% ch3_e (书页 143-150 / p158-p165)
%--D-TAIL--
%JOIN% ，曲率只有一个独立分量。如法炮制，便得到 Ricci 张量
% ---- 书页 143 / p158 ----
\begin{gather*}
R_{xx} = -y^{-2}\\
R_{xy} = 0\\
R_{yy} = -y^{-2},
\tag{3.197}
\end{gather*}

以及标量曲率，
\begin{equation*}
R = -\frac{2}{a^{2}}.
\tag{3.198}
\end{equation*}

我们看到，它与 $S^{2}$ 的情形符号相反，而且特别值得注意的是，它是一个常数。由于我们处在二维，这就足以保证我们的度规确实是最大对称的。当然，还存在另一些坐标，$H^{2}$ 在其中看起来十分不同；习题中会引入其中一种。

于是，在局域上，欧几里得号差的最大对称空间依 $R$ 的符号而为平面、球面或双曲面三者之一。在整体上，任何（欧几里得号差的）最大对称空间都可以这样构造出来：从这三种空间之一中仔细选取一个区域，并把不同的边两两认同，正如平直环面可以由 $\mathbf{R}^{2}$ 构造那样。顺带提一句，让我们简短地谈一谈局域几何与整体拓扑之间的一种联系，它被囊括在 Gauss–Bonnet 定理之中。对于二维紧致无边、可定向的流形，这个定理说
\begin{equation*}
\chi(M) = \frac{1}{4\pi}\int_{M} R\sqrt{|g|}\,\mathrm{d}^{n}x,
\tag{3.199}
\end{equation*}
（译注：原书印作 $\mathrm{d}^{n}x$；就二维曲面而言，按上下文应为 $\mathrm{d}^{2}x$。）

其中 $\chi(M)$ 是该空间的一个拓扑不变量，称为 Euler 示性数（Euler characteristic）。一般说来，它可以由第 2 章提到过的那些上同调空间计算出来；不过在二维情形，它就简单地由下式给出
\begin{equation*}
\chi(M) = 2(1 - g),
\tag{3.200}
\end{equation*}

其中 $g$ 是曲面的亏格（对球面为零，对环面或黎曼面则等于柄的数目）。无论曲率 $R$ 是不是常数，Gauss–Bonnet 定理都成立；然而当 $R$ 是常数时，我们看到所有亏格 $g \geq 2$ 的黎曼面必定具有负曲率，正如球面必定正弯曲、而环面必定平直一样。

继续我们的题外话，暂且来想一想弦论（string theory），它宣称构成宇宙的基本对象是细小的一维弦圈。这样的弦具有二维的“世界面”（world-sheet），而不是一维的世界线。在弦论中做微扰论（相当于量子场论中计算 Feynman 图）涉及对所有世界面几何求和（出于技术上的原因，一般取欧几里得几何）。这听起来像是数目庞大的一堆几何，但在二维中，任何度规都可以写成某个基准度规（fiducial metric）乘以一个共形因子（conformal factor）。这是可信的，因为曲率只有一个分量；习题会请你证明这一点。基准度规可以对每一种世界面拓扑分别选取，而我们不妨把它选成（局域上的）最大对称度规，好让日子好过一些——亏格零用标准的圆球面，亏格一用平面，更高的亏格用双曲面。更加幸运的是，物理上最令人感兴趣的弦论是所谓的临界弦论（critical string theories），对于它们，共形因子本身根本无关紧要。这正是使微扰弦论中的计算成为可能的原因之一；我们只需对一个离散的拓扑集合求和，而且每种拓扑只有有限多个模参数（modular parameters），比如那些告诉我们环面各方向尺寸的参数。

% ---- 书页 144 / p159 ----
我们用最后一个论点来为本节收尾。我们已经探索了具有欧几里得号差的最大对称空间；当然，也存在具有洛伦兹号差的相应时空。我们知道，$R = 0$ 的最大对称时空正是 Minkowski 空间。正曲率的最大对称时空被称为 de Sitter 空间，而负曲率的那个则被颇具想象力地贴上了反 de Sitter 空间的标签。这些时空将在第 8 章中得到更透彻的讨论。

到现在你们应该已经清楚，各附录以重要的方式充实了这些思想。性急的读者可以跳过它们，但那样做未免可惜。

\section{测地偏离}

Riemann 张量作为曲率的后果而露面的又一种方式是：测地偏离。你们无疑听说过，欧几里得（平直）几何的定义性性质是平行公设（parallel postulate）：起初平行的直线永远保持平行。当然，在弯曲空间里这并不成立；在球面上，起初平行的测地线最终必定相交。我们想在任意弯曲空间中把这种行为定量地刻画出来。

问题在于，“平行”这个概念并不能自然地从平直空间延伸到弯曲空间。我们所能做到的最好的事情，是考虑那些可能起初平行的测地线曲线，看看当我们沿着测地线走下去时它们表现如何。为此，我们考虑一个单参数测地线族 $\gamma_{s}(t)$。也就是说，对每个 $s\in\mathbf{R}$，$\gamma_{s}$ 是一条以仿射参量 $t$ 参数化的测地线。这些曲线的集合定义了一张光滑的二维曲面（嵌入在任意维数的流形 $M$ 中）。只要所选的这族测地线互不相交，曲面上的坐标就可以取为 $s$ 和 $t$。整张曲面就是点集 $x^{\mu}(s,t)\in M$。我们有两个自然的矢量场：测地线的切矢量，
\begin{equation*}
T^{\mu} = \frac{\partial x^{\mu}}{\partial t},
\tag{3.201}
\end{equation*}

% ---- 书页 145 / p160 ----
%FIG{3.9}: {一族测地线 $\gamma_{s}(t)$，及其切矢量 $T^{\mu}$。矢量场 $S^{\mu}$ 量度邻近测地线之间的偏离。}

以及偏离矢量（deviation vectors）
\begin{equation*}
S^{\mu} = \frac{\partial x^{\mu}}{\partial s}.
\tag{3.202}
\end{equation*}

这个名称来自如下非正式的观念：$S^{\mu}$ 从一条测地线指向邻近的那些测地线。

$S^{\mu}$ 从一条测地线指向下一条测地线的想法，启发我们定义“测地线的相对速度”（relative velocity of geodesics），
\begin{equation*}
V^{\mu} = (\nabla_{T}S)^{\mu} = T^{\rho}\nabla_{\rho}S^{\mu},
\tag{3.203}
\end{equation*}

以及“测地线的相对加速度”（relative acceleration of geodesics），
\begin{equation*}
A^{\mu} = (\nabla_{T}V)^{\mu} = T^{\rho}\nabla_{\rho}V^{\mu}.
\tag{3.204}
\end{equation*}

这些名字你们听听就好，不必太较真，但这些矢量无疑是良定义的。测地线之间这种\emph{相对}加速度的概念，应当与一条路径偏离测地线运动的加速度区分开来；后者（当 $t$ 为固有时）由 $a^{\mu} = T^{\sigma}\nabla_{\sigma}T^{\mu}$ 给出。

由于 $S$ 和 $T$ 是适配于某个坐标系的基矢量，它们的对易子为零：
\begin{equation*}
[S, T] = 0.
\tag{3.205}
\end{equation*}

由式 (3.37) 我们于是有
\begin{equation*}
S^{\rho}\nabla_{\rho}T^{\mu} = T^{\rho}\nabla_{\rho}S^{\mu}.
\tag{3.206}
\end{equation*}

% ---- 书页 146 / p161 ----
记住这一点，我们来计算加速度：
\begin{align*}
A^{\mu} &= T^{\rho}\nabla_{\rho}(T^{\sigma}\nabla_{\sigma}S^{\mu})\\
&= T^{\rho}\nabla_{\rho}(S^{\sigma}\nabla_{\sigma}T^{\mu})\\
&= (T^{\rho}\nabla_{\rho}S^{\sigma})(\nabla_{\sigma}T^{\mu}) + T^{\rho}S^{\sigma}\nabla_{\rho}\nabla_{\sigma}T^{\mu}\\
&= (S^{\rho}\nabla_{\rho}T^{\sigma})(\nabla_{\sigma}T^{\mu}) + T^{\rho}S^{\sigma}(\nabla_{\sigma}\nabla_{\rho}T^{\mu} + R^{\mu}{}_{\nu\rho\sigma}T^{\nu})\\
&= (S^{\rho}\nabla_{\rho}T^{\sigma})(\nabla_{\sigma}T^{\mu}) + S^{\sigma}\nabla_{\sigma}(T^{\rho}\nabla_{\rho}T^{\mu}) - (S^{\sigma}\nabla_{\sigma}T^{\rho})\nabla_{\rho}T^{\mu}\\
&\qquad + R^{\mu}{}_{\nu\rho\sigma}T^{\nu}T^{\rho}S^{\sigma}\\
&= R^{\mu}{}_{\nu\rho\sigma}T^{\nu}T^{\rho}S^{\sigma}.
\tag{3.207}
\end{align*}

让我们逐行琢磨一下。第一行是 $A^{\mu}$ 的定义；第二行直接来自式 (3.206)；第三行不过是 Leibniz 律；第四行用相反次序的导数加上 Riemann 张量，取代了那个双重协变导数；在第五行我们再次使用 Leibniz 律（这一次与通常的次序相反），然后消去两个相同的项，并注意到含 $T^{\rho}\nabla_{\rho}T^{\mu}$ 的项为零，因为 $T^{\mu}$ 是测地线的切矢量。得到的结果，
\begin{equation*}
\boxed{\,A^{\mu} = \frac{D^{2}}{dt^{2}}S^{\mu} = R^{\mu}{}_{\nu\rho\sigma}T^{\nu}T^{\rho}S^{\sigma}\,,}
\tag{3.208}
\end{equation*}

就是\textbf{测地偏离方程}。它表达了一件我们可能早已预料到的事情：两条相邻测地线之间的相对加速度正比于曲率。

测地偏离方程刻画的是一个单参数相邻测地线族的行为。我们有时也会有兴趣追踪一个多维的相邻测地线集合的行为，比如它可能代表一束光子，或者一份有质量检验粒子的分布。这样的测地线集合构成一个线汇（congruence）；在附录 F 中，我们将推导描述这类线汇演化的方程。

当然，从物理上看，相邻测地线的加速度被解释为引力潮汐力的一种表现。在下一章中，我们将更详细地探讨弯曲时空的性质如何反映在引力场中的物理里。

\section{习题}

% 习题编号照原书
\textbf{1.}\quad 验证度规相容性（$\nabla_{\sigma}g_{\mu\nu}=0$）的下列推论：
\begin{gather*}
\nabla_{\sigma}g^{\mu\nu} = 0\\
\nabla_{\lambda}\epsilon_{\mu\nu\rho\sigma} = 0.
\tag{3.209}
\end{gather*}

% ---- 书页 147 / p162 ----
\textbf{2.}\quad 你们已经熟悉三维欧几里得空间中普通矢量分析里的梯度（$\nabla\phi$）、散度（$\nabla\cdot\mathbf{V}$）和旋度（$\nabla\times\mathbf{V}$）这几种运算。请使用协变导数，在由
\begin{equation*}
x = r\sin\theta\cos\phi
\tag{3.210}
\end{equation*}
\begin{equation*}
y = r\sin\theta\sin\phi
\tag{3.211}
\end{equation*}
\begin{equation*}
z = r\cos\theta.
\tag{3.212}
\end{equation*}
所定义的球坐标 $\{r,\theta,\phi\}$ 中，推导这些运算的公式。把你的结果与 Jackson (1999) 或同类的教材相比较。它们完全相同吗？它们应当完全相同吗？

\textbf{3.}\quad 想象我们有一个\emph{对角}的度规 $g_{\mu\nu}$。证明 Christoffel 符号由下式给出：
\begin{equation*}
\Gamma^{\lambda}_{\mu\nu} = 0
\tag{3.213}
\end{equation*}
\begin{equation*}
\Gamma^{\lambda}_{\mu\mu} = -\frac{1}{2}(g_{\lambda\lambda})^{-1}\partial_{\lambda}g_{\mu\mu}
\tag{3.214}
\end{equation*}
\begin{equation*}
\Gamma^{\lambda}_{\mu\lambda} = \partial_{\mu}\left(\ln\sqrt{|g_{\lambda\lambda}|}\right)
\tag{3.215}
\end{equation*}
\begin{equation*}
\Gamma^{\lambda}_{\lambda\lambda} = \partial_{\lambda}\left(\ln\sqrt{|g_{\lambda\lambda}|}\right)
\tag{3.216}
\end{equation*}
在这些式子中，$\mu\neq\nu\neq\lambda$，而且重复指标\emph{不}作求和。

\textbf{4.}\quad 在欧几里得三维空间中，我们可以经由
\begin{equation*}
x = uv\cos\phi\qquad y = uv\sin\phi\qquad z = \frac{1}{2}(u^{2}-v^{2}).
\end{equation*}
定义抛物面坐标（paraboloidal coordinates）$(u,v,\phi)$。

\textbf{(a)}\quad 求抛物面坐标与笛卡尔坐标之间的坐标变换矩阵 $\partial x^{\alpha}/\partial x^{\beta'}$ 以及逆变换。这个映射中存在奇点吗？

\textbf{(b)}\quad 用笛卡尔的基矢量与基 1-形式表出抛物面坐标的基矢量与基 1-形式。

\textbf{(c)}\quad 求抛物面坐标中的度规与逆度规。

\textbf{(d)}\quad 计算 Christoffel 符号。

\textbf{(e)}\quad 计算散度 $\nabla_{\mu}V^{\mu}$ 与拉普拉斯算符（Laplacian）$\nabla_{\mu}\nabla^{\mu}f$。

\textbf{5.}\quad 考虑一个二维球面，坐标为 $(\theta,\phi)$，度规为
\begin{equation*}
\mathrm{d}s^{2} = \mathrm{d}\theta^{2} + \sin^{2}\theta\,\mathrm{d}\phi^{2}.
\tag{3.217}
\end{equation*}

\textbf{(a)}\quad 证明经线（$\phi=$ 常数）都是测地线，并证明纬线（$\theta=$ 常数）中唯一的一条测地线是赤道（$\theta=\pi/2$）。

\textbf{(b)}\quad 取一个分量为 $V^{\mu}=(1,0)$ 的矢量，让它沿一条等纬度的圆周平行移动一整周。所得到的矢量分量作为 $\theta$ 的函数是什么？

\textbf{6.}\quad 地球表面之外的度规的一个良好近似由下式给出：
\begin{equation*}
\mathrm{d}s^{2} = -(1+2\Phi)\mathrm{d}t^{2} + (1-2\Phi)\mathrm{d}r^{2} + r^{2}(\mathrm{d}\theta^{2} + \sin^{2}\theta\,\mathrm{d}\phi^{2}),
\tag{3.218}
\end{equation*}

% ---- 书页 148 / p163 ----
其中
\begin{equation*}
\Phi = -\frac{GM}{r}
\tag{3.219}
\end{equation*}
可以认为是我们所熟悉的 Newton 引力势。这里 $G$ 是 Newton 引力常数，$M$ 是地球的质量。在本题中可以假设 $\Phi$ 为小量。

\textbf{(a)}\quad 想象在地球表面上、距地心 $R_{1}$ 处有一只钟，而在一座高楼上、距地心 $R_{2}$ 处有另一只钟。试计算每只钟上经历的时间作为坐标时 $t$ 的函数。哪只钟走得更快？

\textbf{(b)}\quad 求出与绕地球赤道（$\theta=\pi/2$）的圆轨道相对应的测地线。$\mathrm{d}\phi/\mathrm{d}t$ 等于多少？

\textbf{(c)}\quad 一颗半径为 $R_{1}$ 的卫星（贴着地面掠过，忽略空气阻力）绕行一整周期间，经历多少固有时？你不妨只算到 $\Phi$ 的一阶。代入地球半径等的实际数值（别忘了把光速恢复进去），得到一个以秒作单位的答案。这个数与静止在地球表面的那只钟上经历的固有时相比如何？

\textbf{7.}\quad 本题中你要利用附录 I 中引入的平行移动传播子（parallel propagator），来看看 Riemann 张量如何从绕一个无穷小回路的平行移动中产生。考虑如下回路：

%FIG{3.11-ex7}: {习题 7 的无穷小回路：由 $x^{1}=0$ 与 $x^{1}=\delta a$、$x^{2}=0$ 与 $x^{2}=\delta b$ 四条曲线围成，顶点为 $A$、$B$、$C$、$D$，箭头标出沿 $A\to B\to C\to D\to A$ 的平行移动方向。}

利用平行移动传播子的无穷级数表达式，计算一个矢量沿此回路从 $A$ 到 $B$ 到 $C$ 到 $D$ 再回到 $A$ 平行移动一周所诱导的变换，准确到 $\delta a$ 和 $\delta b$ 的最低非平凡阶，并证明它正比于 Riemann 张量的相应分量。为了让事情容易些，你可以在旅程的适当路段上取 $x^{1}$ 和 $x^{2}$ 作为参量。

\textbf{8.}\quad 三维球面在坐标 $x^{\mu}=(\psi,\theta,\phi)$ 下的度规可以写成
\begin{equation*}
\mathrm{d}s^{2} = \mathrm{d}\psi^{2} + \sin^{2}\psi(\mathrm{d}\theta^{2} + \sin^{2}\theta\,\mathrm{d}\phi^{2}).
\tag{3.220}
\end{equation*}

% ---- 书页 149 / p164 ----
\textbf{(a)}\quad 计算 Christoffel 联络系数。用什么方法都可以，但通过对积分 (3.49) 作变分来得到联络系数是一种很好的练习。

\textbf{(b)}\quad 计算 Riemann 张量、Ricci 张量和 Ricci 标量。

\textbf{(c)}\quad 证明该度规满足式 (3.191)，从而确认三维球面是一个最大对称空间（正如你们所预期的）。

\textbf{9.}\quad 证明 Weyl 张量 $C^{\mu}{}_{\nu\rho\sigma}$ 在共形变换下保持不变。

\textbf{10.}\quad 证明：对 $n\geq 4$，Weyl 张量满足一种版本的 Bianchi 恒等式，
\begin{equation*}
\nabla_{\rho}C^{\rho}{}_{\sigma\mu\nu} = 2\frac{(n-3)}{(n-2)}\left(\nabla_{[\mu}R_{\nu]\sigma} + \frac{1}{2(n-1)}g_{\sigma[\mu}\nabla_{\nu]}R\right).
\tag{3.221}
\end{equation*}

\textbf{11.}\quad 既然具有度规 (3.192) 的庞加莱半平面是最大对称的，我们或许会预期它绕任何一点都具有旋转对称性，尽管这在 $\{x,y\}$ 坐标中当然很不明显。如果确实如此，那就应当能把度规写成旋转对称性一目了然的形式，例如
\begin{equation*}
\mathrm{d}s^{2} = f^{2}(r)[\mathrm{d}r^{2} + r^{2}\mathrm{d}\theta^{2}].
\tag{3.222}
\end{equation*}
为了证明这样做行得通，请计算这个度规的标量曲率，并在处处 $R=-2/a^{2}$ 的条件下解出函数 $f(r)$。坐标 $r$ 的取值范围是什么？

\textbf{12.}\quad 证明任何 Killing 矢量 $K^{\mu}$ 都满足正文中提到的下列关系式：
\begin{gather*}
\nabla_{\mu}\nabla_{\sigma}K^{\rho} = R^{\rho}{}_{\sigma\mu\nu}K^{\nu}\\
K^{\lambda}\nabla_{\lambda}R = 0.
\tag{3.223}
\end{gather*}

\textbf{13.}\quad 对下列空间，求出完备 Killing 矢量场集合的显式表达式：

\textbf{(a)}\quad Minkowski 空间，度规为 $\mathrm{d}s^{2}=-\mathrm{d}t^{2}+\mathrm{d}x^{2}+\mathrm{d}y^{2}+\mathrm{d}z^{2}$。

\textbf{(b)}\quad 一个具有坐标 $\{u,v,x,y\}$ 和度规
\begin{equation*}
\mathrm{d}s^{2} = -(\mathrm{d}u\,\mathrm{d}v + \mathrm{d}v\,\mathrm{d}u) + a^{2}(u)\mathrm{d}x^{2} + b^{2}(u)\mathrm{d}y^{2},
\tag{3.224}
\end{equation*}
的时空，其中 $a$ 与 $b$ 是 $u$ 的未指定函数。这代表一个引力波时空。（\emph{提示}，不必给出证明：一共有五个 Killing 矢量，而且它们全都有 $u$ 分量 $K^{u}$ 为零。）

在所有这些情形中，都要当心上指标与下指标的区别。

\textbf{14.}\quad 考虑二维球面的那三个 Killing 矢量，即式 (3.188)。证明它们的对易子满足如下代数关系：
\begin{gather*}
[R, S] = T\\
[S, T] = R\\
[T, R] = S.
\tag{3.225}
\end{gather*}

\textbf{15.}\quad 利用附录 F 中讨论的 Raychaudhuri 方程证明：如果一种流体沿测地线流过时空，切变与膨胀均为零，那么该时空必定拥有一个类时 Killing 矢量。

% ---- 书页 150 / p165 ----
\textbf{16.}\quad 再次考虑三维球面上的度规，
\begin{equation*}
\mathrm{d}s^{2} = \mathrm{d}\psi^{2} + \sin^{2}\psi(\mathrm{d}\theta^{2} + \sin^{2}\theta\,\mathrm{d}\phi^{2}).
\tag{3.226}
\end{equation*}
本题我们要用到附录 J 中讨论的非坐标基。在一组正交归一的 1-形式标架 $\hat{\theta}^{(a)}$ 下，度规将变成
\begin{equation*}
\mathrm{d}s^{2} = \hat{\theta}^{(1)}\otimes\hat{\theta}^{(1)} + \hat{\theta}^{(2)}\otimes\hat{\theta}^{(2)} + \hat{\theta}^{(3)}\otimes\hat{\theta}^{(3)}.
\tag{3.227}
\end{equation*}

\textbf{(a)}\quad 找出这样一组正交归一的 1-形式标架，使得矩阵 $e^{\mu}{}_{a}$ 是对角的。不必担心覆盖整个流形。

\textbf{(b)}\quad 通过求解 $\mathrm{d}e^{a} + \omega^{a}{}_{b}\wedge e^{b} = 0$，计算自旋联络（spin connection）的分量。

\textbf{(c)}\quad 先计算曲率 2-形式 $R^{a}{}_{b\mu\nu}$ 的分量，再加以转换，从而算出适配于 $x^{\mu}$ 的坐标基中 Riemann 张量 $R^{\rho}{}_{\sigma\mu\nu}$ 的分量。
